(1) 理想谱域OCT中，一面弱反射镜的参考臂与样品臂产生谱干涉 \(I(k)=S(k)[1+V\cos(2kz_0+\phi_0)]\)，\(S(k)=S_0e^{-(k-k_0)^2/(2\sigma_k^2)}/(\sqrt{2\pi}\sigma_k)\)，其中 \(k>0\) 近似可延伸到整实轴且 \(k_0\gg\sigma_k\)。求背景扣除后的复傅里叶深度响应。
(2) 由第一问推导轴向复振幅包络的半高全宽（FWHM），并分别写出强度 \(|A|^2\) 的FWHM，解释为何两者相差 \(\sqrt2\)。
(3) 若干涉谱中额外出现相位 \(\phi_2(k-k_0)^2/2\)，推导重建深度包络的展宽因子和峰值幅度衰减因子，给出理想数值色散补偿。
(4) 谱仪以等间距 \(\Delta k\) 采样 \(k\)，说明傅里叶深度的混叠周期和可无混叠的有符号深度区间；分析 \(V\to0\) 与 \(\sigma_k\to0\) 两个物理边界。
